% Part: computability
% Chapter: recursive-functions
% Section: general-recursive-functions

\documentclass[../../../include/open-logic-section]{subfiles}

\begin{document}

\olfileid[pt]{cmp}{rec}{gen}
\olsection{Funções recursivas gerais}

Há outra maneira de obter um conjunto de funções totais. Dizemos que
uma função total $f(x,\vec z)$ é \emph{regular} se, para toda sequência
de números naturais $\vec z$, existe um $x$ tal que $f(x,\vec z) = 0$.
Em outras palavras, as funções regulares são exatamente aquelas às
quais se pode aplicar a busca ilimitada e obter como resultado uma
função total. Podemos, de modo conservador, restringir a busca
ilimitada às funções regulares:

\begin{defn}
\ollabel{defn:general-recursive}
O conjunto das \emph{funções recursivas gerais} é o menor conjunto de
funções dos números naturais nos números naturais (de diversas
aridades) que contém zero, sucessor e projeções e é fechado sob
composição, recursão primitiva e busca ilimitada aplicada a funções
\emph{regulares}.
\end{defn}

É evidente que toda função recursiva geral é total. A diferença entre
\olref{defn:general-recursive} e
\olref[par]{defn:recursive-fn} é que, na última, permite-se usar funções
recursivas parciais durante o processo; a única exigência é que a
função obtida ao final seja total. Assim, a palavra ``geral'', um
vestígio histórico, é uma denominação inadequada; à primeira vista,
\olref{defn:general-recursive} é \emph{menos} geral do que
\olref[par]{defn:recursive-fn}. Felizmente, porém, a diferença é
ilusória; embora as definições sejam diferentes, o conjunto das
funções recursivas gerais e o conjunto das funções recursivas são um
só e o mesmo.

\end{document}
